\section{The full Hecke algebra and the theorem}
\label{sec:introduction}

Congruences between modular forms of different weights assemble their
Hecke eigenvalues into a single $p$-adic algebra. Its points need not
come from classical forms of one weight. The resulting modularity
question asks whether a Galois representation occurs in this larger
spectrum, even when no Hodge-theoretic condition is imposed at $p$.
We prove this occurrence statement for every odd, absolutely
irreducible two-dimensional representation over $\Q$ at $p=2$.

Fix an odd positive integer $N$. For $i\geq1$, let
$M_i(\Gamma_1(N))$ denote the space of all complex holomorphic modular
forms of weight $i$, including Eisenstein forms. Let
$T_{\leq k}^{(2)}(N)$ be the $\Z$-algebra acting on
$\bigoplus_{i=1}^{k}M_i(\Gamma_1(N))$
generated by $T_\ell$ and $\ell S_\ell$ for primes $\ell\nmid2N$.
On weight $i$, the operator $S_\ell$ is
$\ell^{i-2}\langle\ell\rangle$. Set
\begin{equation}\label{eq:full-hecke}
 \mathbb T_2(N)=
 \varprojlim_k\bigl(\Z_2\otimes_{\Z}T_{\leq k}^{(2)}(N)\bigr).
\end{equation}
Each finite-stage algebra has its $2$-adic topology, and
$\mathbb T_2(N)$ has the resulting inverse-limit topology.
We fix the Frobenius convention in which a classical eigenform has
characteristic polynomial
\begin{equation}\label{eq:frobenius}
 Z^2-T_\ell Z+\ell S_\ell
\end{equation}
at a good prime, as in \cite[Equation~(2.2)]{Dim}.

\begin{theorem}\label{thm:main}
Let $E/\Q_2$ be finite, with ring of integers $\OO_E$, and let
$r:G_{\Q}\longrightarrow\GL_2(E)$
be continuous and absolutely irreducible. Suppose that $r$ is
unramified outside finitely many finite primes and is odd:
$\det r(c)=-1$ for complex conjugation $c$.
Then there exist an odd positive integer $N$, divisible by every odd
prime where $r$ ramifies, and a continuous $\Z_2$-algebra homomorphism
\[
 \lambda:\mathbb T_2(N)\longrightarrow\OO_E
\]
such that, for every prime $\ell\nmid2N$,
\[
 \lambda(T_\ell)=\tr r(\operatorname{Frob}_\ell),\qquad
 \lambda(\ell S_\ell)=\det r(\operatorname{Frob}_\ell).
\]
\end{theorem}

We call this conclusion \emph{pro-modularity}. No condition on the
semisimplified reduction of $r$ is imposed. In particular, the theorem
includes scalar residual representations. The compact image of $r$
preserves an $\OO_E$-lattice, so its traces and determinants are integral.

\subsection*{Context and main ideas}

The Fontaine--Mazur conjecture relates Galois representations satisfying
geometric local conditions to algebraic geometry.  In its odd,
two-dimensional regular form over $\Q$, the expected conclusion is
classical cuspidal modularity up to a Tate twist
\cite{FM1995,Pan2022}.  Removing the local condition leads to a different
question: occurrence of the representation in a completed Hecke algebra.
Emerton records this expectation for the full varying-weight algebra in
\cite[Conjecture~2.12]{Emerton2011Families}.  That formulation prescribes
the allowed tame primes.  Theorem~\ref{thm:main} proves dyadic occurrence
for absolutely irreducible representations while allowing auxiliary
primes in the tame level.

Skinner and Wiles developed pro-modularity arguments through ordinary
families in the residually reducible setting
\cite[Introduction and Section~4.1]{SkinnerWiles1999}.
Emerton's local--global compatibility work separates pro-modularity from
the production of classical vectors
\cite[Section~1.2]{Emerton2011}; the pro-modularity theorem recorded there
assumes $p>2$, residual irreducibility, and local residual restrictions.
Pan extends the family method to the residually reducible nonordinary
setting at odd primes: ordinary loci supply modular points, and patching
at one-dimensional primes propagates Hecke support
\cite[Introduction and Theorem~4.1.7]{Pan2022}.
His classical modularity theorems retain regular de Rham hypotheses.
These precedents suggest placing an unrestricted representation in a
family and finding a useful geometric point elsewhere in that family.

The local representation theory used here comes from Colmez's
$p$-adic Langlands correspondence \cite{Colmez2010}, in the all-prime form
of Colmez, Dospinescu, and Pa\v{s}k\=unas
\cite[Theorem~1.1]{CDP2014}.  Pa\v{s}k\=unas and Tung supply the integral
block finiteness and Cayley--Hamilton descriptions needed at $2$,
including scalar residual blocks
\cite[Theorems~1.2--1.4 and Section~4.1]{PaskunasTung2021}.
Their results control local parameters but do not by themselves place a
global representation on Hecke support.

After a continuous cyclotomic twist, we place the original
representation in a fixed-determinant trace family $D$ of dimension
at least three. The curve and connectedness arguments of \cite{FM2}
put this family on Hecke support after solvable totally real base
change, as developed in Section~\ref{sec:family}. The main task is
then to find a regular de Rham point on $D$, from which Hecke
support can be transferred to the whole family over $\Q$.

Two constructions proved here produce that point. We choose the
base-change field to split completely at $2$, so that all dyadic
restrictions of the family have the same local Galois parameter.
Section~\ref{sec:diagonal} passes through the local block equivalences
to construct an object for one copy of $\GL_2(\Q_2)$ from an admissible
representation of the product of dyadic groups. Equality of the parameters
ensures that a failure of
admissibility would persist on a characteristic-two curve in every factor,
contradicting the admissibility of the product representation.
This comparison includes scalar residual blocks.
Section~\ref{sec:regular-point} then uses a chain of prime
specializations in $D$ to force maximal growth, and hence positive
Iwasawa rank, for the one-factor object. Polynomial approximation
detects a locally algebraic depth-zero supercuspidal type, yielding
a regular de Rham point $x$. The regular Fontaine--Mazur theorem
\cite[Theorem~1.1]{FM2} makes $x$ classical up to Tate twist, and
Section~\ref{sec:hecke} puts its classical packet and continuous
cyclotomic twists in the precise algebra \eqref{eq:full-hecke}.

The final passage uses the dimension theorem of
\cite[Theorem~1.1]{Dim} and the adjoint Selmer vanishing theorem of
Newton and Thorne \cite[Theorem~5.4]{NewtonThorne2023}.
At $x$, the global deformation problem with varying determinant has
tangent dimension at most three, while a dimension-four integral
Hecke component has local dimension three. Equality makes the
ambient local ring regular and forces its Hecke kernel to vanish
there. Since $D$ is a domain through $x$, the same kernel vanishes
on all of $D$. Section~\ref{sec:conclusion} carries out this comparison
and recovers the original coefficient point. The Hodge-theoretic
condition is needed only at the auxiliary point.

\subsection*{Coefficient and deformation conventions}

We allow finite extensions of coefficient fields during the proof.
Write $\OO$ for the current ring of integers, $\varpi$ for a
uniformizer, and $\kk$ for the residue field. A \emph{coefficient
point} of a complete local $\OO$-algebra is a continuous map to the
ring of integers of a finite coefficient extension. Its kernel is
a characteristic-zero prime; the corresponding integral quotient
has dimension one. A domain over $\OO$ is \emph{horizontal} if
$\varpi$ is nonzero in it.

At residue characteristic two, a pseudorepresentation means a
two-dimensional determinant law, with both trace and determinant
data, in the sense of \cite{Chenevier2014}. Thus no division by $2$
is implicit in an integral deformation ring. All global deformation
problems have a fixed finite ramification set containing $2$ and
allow ramification at infinity. Trace images and support containments
are closed images and set-theoretic containments; the domains used
for the family are reduced.

For the local categories and completed definite forms we use the
normalization of \cite[Section~2]{FM2}. If the fixed global
determinant is $\chi$, the central character is $\psi=\chi\eps$,
where $\eps$ is the cyclotomic character. The local block parameter
is the Galois parameter twisted by $\eps$, of determinant
$\psi\eps=\chi\eps^2$. Group duals have the contragredient action;
commuting coefficient and Hecke operators act by transposition.
The conversion to the Frobenius identity \eqref{eq:frobenius} is
made explicitly when passing to classical Hecke points.
